% math0 manuscript. Build with: latexmk -pdf paper.tex
\documentclass[11pt]{amsart}
\usepackage[margin=1.1in]{geometry}
\usepackage{amsmath,amssymb,amsthm,mathtools}
\usepackage{url}
\usepackage{hyperref}
\usepackage{xcolor}
\usepackage{listings}
\usepackage{newunicodechar}

\lstdefinelanguage{lean}{
  keywords={theorem,lemma,def,abbrev,example,fun,by,have,show,from,let,in,if,then,else,match,with,
    namespace,end,open,import,section,variable,instance,structure,where,calc,obtain,intro,exact,
    refine,rcases,rintro,simp,omega,decide,norm_num,linarith,nlinarith,positivity,induction,cases,
    private,noncomputable,rw,constructor,by_contra},
  sensitive=true, comment=[l]{--}, morecomment=[s]{/-}{-/}, string=[b]"}
\lstset{language=lean, basicstyle=\ttfamily\footnotesize, columns=fullflexible, keepspaces=true, breaklines=true,
  keywordstyle=\color{blue!60!black}\bfseries, commentstyle=\color{gray}, frame=single,
  extendedchars=true, inputencoding=utf8,
  literate={ℕ}{{$\mathbb{N}$}}1 {∀}{{$\forall$}}1 {∃}{{$\exists$}}1 {∈}{{$\in$}}1
    {∧}{{$\wedge$}}1 {∨}{{$\vee$}}1 {¬}{{$\neg$}}1 {→}{{$\to$}}1 {←}{{$\leftarrow$}}1
    {↔}{{$\leftrightarrow$}}1 {≤}{{$\le$}}1 {≡}{{$\equiv$}}1 {∣}{{$\mid$}}1
    {⟨}{{$\langle$}}1 {⟩}{{$\rangle$}}1 {·}{{$\cdot$}}1}
% Unicode commonly found in Lean sources.
\newunicodechar{ℕ}{\ensuremath{\mathbb{N}}}
\newunicodechar{ℤ}{\ensuremath{\mathbb{Z}}}
\newunicodechar{ℚ}{\ensuremath{\mathbb{Q}}}
\newunicodechar{ℝ}{\ensuremath{\mathbb{R}}}
\newunicodechar{ℂ}{\ensuremath{\mathbb{C}}}
\newunicodechar{∀}{\ensuremath{\forall}}
\newunicodechar{∃}{\ensuremath{\exists}}
\newunicodechar{→}{\ensuremath{\to}}
\newunicodechar{←}{\ensuremath{\leftarrow}}
\newunicodechar{↔}{\ensuremath{\leftrightarrow}}
\newunicodechar{↦}{\ensuremath{\mapsto}}
\newunicodechar{¬}{\ensuremath{\neg}}
\newunicodechar{∧}{\ensuremath{\wedge}}
\newunicodechar{∨}{\ensuremath{\vee}}
\newunicodechar{≤}{\ensuremath{\le}}
\newunicodechar{≥}{\ensuremath{\ge}}
\newunicodechar{≠}{\ensuremath{\ne}}
\newunicodechar{∈}{\ensuremath{\in}}
\newunicodechar{∉}{\ensuremath{\notin}}
\newunicodechar{⊆}{\ensuremath{\subseteq}}
\newunicodechar{∩}{\ensuremath{\cap}}
\newunicodechar{∪}{\ensuremath{\cup}}
\newunicodechar{∑}{\ensuremath{\sum}}
\newunicodechar{∏}{\ensuremath{\prod}}
\newunicodechar{⟨}{\ensuremath{\langle}}
\newunicodechar{⟩}{\ensuremath{\rangle}}
\newunicodechar{·}{\ensuremath{\cdot}}
\newunicodechar{∘}{\ensuremath{\circ}}
\newunicodechar{⁻}{\ensuremath{^{-}}}
\newunicodechar{¹}{\ensuremath{^{1}}}
\newunicodechar{²}{\ensuremath{^{2}}}
\newunicodechar{₀}{\ensuremath{_{0}}}
\newunicodechar{₁}{\ensuremath{_{1}}}
\newunicodechar{₂}{\ensuremath{_{2}}}
\newunicodechar{α}{\ensuremath{\alpha}}
\newunicodechar{β}{\ensuremath{\beta}}
\newunicodechar{ε}{\ensuremath{\varepsilon}}
\newunicodechar{δ}{\ensuremath{\delta}}
\newunicodechar{π}{\ensuremath{\pi}}
\newunicodechar{σ}{\ensuremath{\sigma}}
\newunicodechar{φ}{\ensuremath{\varphi}}
\newunicodechar{∞}{\ensuremath{\infty}}
\newunicodechar{⌊}{\ensuremath{\lfloor}}
\newunicodechar{⌋}{\ensuremath{\rfloor}}
\newunicodechar{⌈}{\ensuremath{\lceil}}
\newunicodechar{⌉}{\ensuremath{\rceil}}
\newunicodechar{∣}{\ensuremath{\mid}}
\newunicodechar{≡}{\ensuremath{\equiv}}
\newunicodechar{×}{\ensuremath{\times}}
\newunicodechar{⧸}{\ensuremath{/}}
\newunicodechar{‖}{\ensuremath{\|}}
\newunicodechar{√}{\ensuremath{\surd}}
\newunicodechar{∅}{\ensuremath{\emptyset}}
\newunicodechar{⦃}{\ensuremath{\{\!|}}
\newunicodechar{⦄}{\ensuremath{|\!\}}}
\newunicodechar{𝓝}{\ensuremath{\mathcal{N}}}
\newunicodechar{≫}{\ensuremath{\gg}}
\newunicodechar{≪}{\ensuremath{\ll}}

\theoremstyle{plain}
\newtheorem{theorem}{Theorem}
\newtheorem{lemma}[theorem]{Lemma}
\newtheorem{proposition}[theorem]{Proposition}
\newtheorem{corollary}[theorem]{Corollary}
\theoremstyle{definition}
\newtheorem{definition}[theorem]{Definition}
\newtheorem{remark}[theorem]{Remark}

\emergencystretch=3em
% Lean identifiers in running text (breakable at dots and underscores).
\expandafter\def\expandafter\UrlBreaks\expandafter{\UrlBreaks\do\_\do\.}
\DeclareUrlCommand\decl{\urlstyle{tt}}
\newcommand{\N}{\mathbb{N}}
\newcommand{\cD}{\mathcal{D}}

\title[Irvine's conjecture for OEIS A062887]{An odd number has a multiple whose decimal digits are odd and decrease by two modulo $10$ if and only if it is not divisible by $125$: Irvine's conjecture for OEIS A062887}
\author{math0}
\address{Produced by Claude (an AI model developed by Anthropic) in the math0 project}
\date{October 19, 2026}

\begin{document}
\begin{abstract}
For $n\ge0$ let $a(n)$ be the smallest positive multiple of $2n+1$ whose decimal digits are odd and, read from left to right, each two less modulo $10$ than the preceding one, or $0$ if there is no such multiple; this is the sequence A062887 of the OEIS, contributed by Amarnath Murthy in 2001.
We prove the conjecture recorded in the entry by Sean A.~Irvine in 2023 that $a(n)=0$ if and only if $125$ divides $2n+1$: the last three digits of such a number are never divisible by $125$, and otherwise a suitable number of repetitions of the block $31975$, chosen with Euler's theorem, is a multiple of $2n+1$.
To our knowledge this is the first proof of the conjecture; the argument is elementary.
The theorem is verified in Lean~4 against a statement that math0 formalized from the OEIS entry, reviewed for fidelity and frozen before the proof attempt, whose definitions are checked by the kernel against all $26$ listed terms; this statement is not part of the formal-conjectures repository.
\end{abstract}
\maketitle

\section{Introduction}

Call a positive integer \emph{admissible} if its decimal digits are odd and, read from left to right, each digit is two less modulo $10$ than the previous one. Reading from the left, the digits of an admissible number run cyclically through $\dots,9,7,5,3,1,9,7,\dots$; examples are $1$, $75$, $319$, $975$ and $9753197$.
The sequence A062887 of the OEIS \cite{A062887}, contributed by Amarnath Murthy on June~28, 2001, with offset $0$, is
\[
  a(n)=\text{the smallest admissible multiple of }2n+1,\ \text{or }0\text{ if there is none}.
\]
The entry lists $a(0),\dots,a(25)$, with more terms by Sascha Kurz (2002); for example $a(5)=319$ and $a(10)=197531975319$, and the largest listed term, $a(25)=753197531975319753$, has $18$ digits.
On April~14, 2023, Sean A.~Irvine added the conjecture that $a(n)=0$ exactly when $125\mid 2n+1$. We prove it (Theorem~\ref{thm:main}).

\subsection*{Method}
Since $125\mid1000$, an admissible number is congruent modulo $125$ to its last three digits, and these are never divisible by $125$. Conversely, if $125\nmid m=2n+1$ and $m'$ is the largest divisor of $m$ coprime to $5$, the admissible number formed by $q=\varphi(99999\,m')$ copies of the block $31975$ is divisible by $25m'$, hence by $m$, by Euler's theorem.

\subsection*{Related sequences}
A062887 cross-references A062885 \cite{A062885}, also due to Murthy: the smallest multiple of $n$ whose digits are even and each two less (mod $10$) than the previous one, or $-1$ if there is none. For it, John W.~Layman showed in 2001 that the value at $n=25$ is $-1$, and the entry links an analysis by Don Reble.
In the sibling sequence A062886 \cite{A062886} the digits are odd and each is two \emph{more} (mod $10$) than the previous one. Nathaniel Johnston observed in 2011 that, for $n\le500$, its zeros are exactly the $n$ with $25\mid 2n+1$, and the entry's example for $a(12)=0$ gives the reason why $25\mid 2n+1$ forces a zero (a multiple of $25$ ends in $00$, $25$, $50$ or $75$). The argument of this note proves Johnston's observation for all $n$ (Remark~\ref{rem:886}).

\subsection*{Status and prior work}
The statement is still marked ``Conjecture'' in A062887; the current revision (\#12, April~14, 2023) dates from the day the conjecture was added. We read the entry in the copy of the OEIS database used by math0, dated October~17, 2026, and no other entry in that copy refers to A062887.
The math0 referees' web searches (including queries aimed at the OEIS, MathOverflow and Mathematics Stack Exchange, and searches for the repeated blocks $31975$ and $97531$ together with divisibility by $125$) found no proof or counterexample.
The statement has no file in the formal-conjectures repository \cite{FC} at the commit pinned by math0, and none among the AlphaProof Nexus outputs \cite{APN}.
We could not read Reble's analysis of A062885 (our copy of the database contains only a Git LFS pointer to it); it concerns the even-digit sequence.
We therefore classify the result as new, with the caveat that the argument is elementary, and that the direction ``$125\mid 2n+1\Rightarrow a(n)=0$'' is the same kind of last-digits argument that the entry A062886 gives for the modulus $25$.

\section{Statement}\label{sec:statement}

Every $N\ge1$ has a unique decimal expansion $N=\sum_{i=0}^{L-1}d_i10^i$ with $d_i\in\{0,\dots,9\}$ and $d_{L-1}\ne0$; we call $(d_0,\dots,d_{L-1})$ its (little-endian) \emph{digit list} and $L$ its number of digits. Read from left to right the digits are $d_{L-1},d_{L-2},\dots,d_0$.

\begin{definition}\label{def:adm}
A positive integer $N$ with digit list $(d_0,\dots,d_{L-1})$ is \emph{admissible} if every $d_i$ is odd and $d_i+2\equiv d_{i+1}\pmod{10}$ for $0\le i<L-1$. Let $\cD$ be the set of admissible numbers, and for $n\ge0$ let $a(n)=\min\{N\in\cD: 2n+1\mid N\}$ if this set is nonempty and $a(n)=0$ otherwise.
\end{definition}

The name of A062887 \cite{A062887} reads, verbatim:
\begin{quote}
{\small\texttt{Smallest multiple of 2n+1 with the property that its digits are odd and each digit is two less (mod 10) than the previous digit, or 0 if no such number exists.}}
\end{quote}
and the COMMENTS section contains, verbatim:
\begin{quote}
{\small\texttt{Conjecture: The only numbers n for which a(n)=0 are those for which 2*n+1 is divisible by 125. - \_Sean A. Irvine\_, Apr 14 2023}}
\end{quote}
Definition~\ref{def:adm} is the name with ``the previous digit'' taken in the left-to-right reading: $d_i$ is two less than $d_{i+1}$ modulo $10$. This agrees with the entry's Maple program, whose candidates have digits $(2j+1-2h)\bmod 10$ for $h=1,2,\dots$ from the left, and with the listed data.\footnote{The entry's example line reads ``\texttt{a(7) = 975 = 13*75}''. With offset $0$, the listed data give $a(6)=975$ (and $2\cdot6+1=13$) and $a(7)=75$.}

\begin{theorem}\label{thm:main}
For every $n\ge0$, $a(n)=0$ if and only if $125\mid 2n+1$. Equivalently, an odd positive integer $m$ has an admissible multiple if and only if $125\nmid m$.
\end{theorem}

Irvine's sentence asserts that $a(n)=0$ implies $125\mid 2n+1$; Theorem~\ref{thm:main} also contains the converse, which is the easy direction (Lemma~\ref{lem:125}).

\subsection*{Lean formalization}
The statement was formalized by math0 in \path{lean/Math0/Statements/OEIS/A062887.lean}, reviewed for fidelity and frozen (SHA-256 \texttt{544c69cd\ldots a828}) before any proof attempt. Its definitions of the sequence and the statement read, verbatim (the import, comments and docstrings, blank lines, the decidability instance, the auxiliary definitions \decl{gen}, \decl{cands} and \decl{aSearch} used for the data check, and all theorems are omitted):
\begin{lstlisting}
namespace Math0.Statements.A062887
def IsDesc2 (N : ℕ) : Prop :=
  0 < N ∧ (∀ d ∈ Nat.digits 10 N, Odd d) ∧
    (Nat.digits 10 N).reverse.IsChain (fun prev next => next + 2 ≡ prev [MOD 10])
open Classical in
noncomputable def a (n : ℕ) : ℕ :=
  if h : ∃ N, IsDesc2 N ∧ 2 * n + 1 ∣ N then Nat.find h else 0
def statement : Prop := ∀ n : ℕ, a n = 0 ↔ 125 ∣ 2 * n + 1
end Math0.Statements.A062887
\end{lstlisting}
Here \texttt{Nat.digits 10 N} is the digit list $(d_0,\dots,d_{L-1})$ defined above (empty for $N=0$), so its reverse lists the digits from left to right, and \texttt{IsChain R} requires $R$ between each pair of consecutive entries. Thus \texttt{IsDesc2 N} says that $N\in\cD$; the condition $0<N$ excludes $N=0$, whose empty digit list would satisfy the other two conditions vacuously. \texttt{Nat.find h} is the least $N$ with \texttt{IsDesc2 N} and $2n+1\mid N$, so \decl{a} is the function $a$ of Definition~\ref{def:adm}. The statement quantifies over all $n\ge0$ without further hypotheses.

To compare \decl{a} with the OEIS data, the statement file describes $\cD$ explicitly and proves the following (statements shown, proofs omitted except for \decl{a_data}):
\begin{lstlisting}
def gen : ℕ → ℕ → List ℕ
  | _, 0 => []
  | x, L + 1 => x :: gen ((x + 2) % 10) L
theorem isDesc2_ofDigits_gen (x L : ℕ) (hx : x < 10) (hodd : Odd x) (hL : 1 ≤ L) :
    IsDesc2 (Nat.ofDigits 10 (gen x L))
theorem exists_gen_of_isDesc2 {N : ℕ} (h : IsDesc2 N) :
    ∃ x L, x < 10 ∧ Odd x ∧ 1 ≤ L ∧ Nat.digits 10 N = gen x L
theorem a_eq_of_aSearch {n B v : ℕ} (h : aSearch n B = some v) : a n = v
theorem a_data : (List.range 26).map (fun i => a (i + 0)) =
    [1, 3, 5, 7, 9, 319, 975, 75, 7531, 19, 197531975319, 3197, 75, 5319, 319, 31,
      3197531975319, 53197531975, 1975319, 975, 531975, 1975319753197, 19753197531975,
      75319753197531, 319753197531975319, 753197531975319753] := by
  simp only [Nat.add_zero]
  exact map_a_eq_of_aSearch (B := 18) (by decide +kernel)
\end{lstlisting}
The list \texttt{gen x L} is $\bigl(x,(x+2)\bmod10,(x+4)\bmod10,\dots\bigr)$ of length $L$, and the first two theorems say that $\cD$ consists exactly of the numbers with digit list \texttt{gen x L}, $x$ an odd digit and $L\ge1$ (Lemma~\ref{lem:gen} below). The computable function \texttt{aSearch n B} (defined in the file through a list \decl{cands} of candidates) returns the least multiple of $2n+1$ among the $5B$ admissible numbers \texttt{Nat.ofDigits 10 (gen x L)} with $x\in\{1,3,5,7,9\}$ and $1\le L\le B$, if there is one. By \decl{a_eq_of_aSearch}, a value \texttt{some v} determines \texttt{a n = v}: an admissible $k<v$ has at most as many digits as $v$, hence lies in the search range, so $k$ is not a multiple of $2n+1$. In \decl{a_data} the kernel evaluates \texttt{aSearch i 18} for $i=0,\dots,25$ and compares the results with the $26$ listed terms; the lemma \decl{map_a_eq_of_aSearch} applies \decl{a_eq_of_aSearch} termwise. The file also checks the example $975=13\cdot75$.

\section{Proof}\label{sec:proof}

We follow the Lean file \path{lean/Math0/Frontier/OEIS/A062887.lean}; each step names the corresponding declaration. For an odd digit $x\in\{1,3,5,7,9\}$ and $L\ge0$ put
\[
  \gamma(x,L)=\bigl((x+2i)\bmod 10\bigr)_{0\le i<L},\qquad G(x,L)=\sum_{i=0}^{L-1}\bigl((x+2i)\bmod 10\bigr)\,10^i ,
\]
so $\gamma(x,L)$ is \texttt{gen x L} and $G(x,L)$ is \texttt{Nat.ofDigits 10 (gen x L)}. Since $(x+2)\bmod10$ is again an odd digit,
\begin{equation}\label{eq:rec}
  \gamma(x,L+1)=\bigl(x,\gamma((x+2)\bmod 10,L)\bigr),\qquad G(x,L+1)=x+10\,G\bigl((x+2)\bmod 10,L\bigr).
\end{equation}

\begin{lemma}[\decl{isDesc2_ofDigits_gen}, \decl{exists_gen_of_isDesc2}]\label{lem:gen}
A positive integer $N$ is admissible if and only if $N=G(x,L)$ for some $x\in\{1,3,5,7,9\}$ and $L\ge1$. In that case $\gamma(x,L)$ is the digit list of $N$.
\end{lemma}

\begin{proof}
All entries of $\gamma(x,L)$ are odd digits; in particular the last one is nonzero, so $\gamma(x,L)$ is the digit list of $G(x,L)$ (Mathlib's \decl{Nat.digits_ofDigits}). Consecutive entries satisfy $d_{i+1}=(d_i+2)\bmod 10\equiv d_i+2$, and $G(x,L)\ge x\ge1$ for $L\ge1$; so $G(x,L)$ is admissible.
Conversely, let $N$ be admissible with digit list $(d_0,\dots,d_{L-1})$; then $L\ge1$ because $N>0$. For each $i$, $d_{i+1}\equiv d_i+2\pmod{10}$ and $0\le d_{i+1}<10$ give $d_{i+1}=(d_i+2)\bmod 10$, so by induction along the list $(d_0,\dots,d_{L-1})=\gamma(d_0,L)$, and $d_0$ is an odd digit. (In Lean, \decl{List.isChain_reverse} turns the condition on the reversed list for $(\mathit{prev},\mathit{next})\mapsto \mathit{next}+2\equiv \mathit{prev}$ into the condition on the list itself for $(p,q)\mapsto p+2\equiv q$, and the induction is \decl{gen_of_chain}.)
\end{proof}

\begin{lemma}[\decl{a_eq_zero_iff}]\label{lem:zero}
For every $n\ge0$, $a(n)=0$ if and only if no admissible number is divisible by $2n+1$.
\end{lemma}

\begin{proof}
If some admissible number is divisible by $2n+1$, then $a(n)$ is the least such number; it is admissible, hence positive. Otherwise $a(n)=0$ by definition.
\end{proof}

\begin{lemma}[\decl{not_dvd_125}]\label{lem:125}
No admissible number is divisible by $125$.
\end{lemma}

\begin{proof}
Let $N$ be admissible. By Lemma~\ref{lem:gen}, $N=G(x,L)$ with $x\in\{1,3,5,7,9\}$ and $L\ge1$. Applying \eqref{eq:rec} up to three times gives the following values, where $K=G\bigl((x+6)\bmod 10,L-3\bigr)\ge0$ for $L\ge3$:
\[
\begin{array}{c|ccccc}
 x & 1 & 3 & 5 & 7 & 9\\\hline
 N\ (L=1) & 1 & 3 & 5 & 7 & 9\\
 N\ (L=2) & 31 & 53 & 75 & 97 & 19\\
 N-1000K\ (L\ge3) & 531 & 753 & 975 & 197 & 319\\
 (N-1000K)\bmod 125 & 31 & 3 & 100 & 72 & 69
\end{array}
\]
If $L\le2$, then $0<N<125$. If $L\ge3$, then $N\equiv N-1000K\not\equiv0\pmod{125}$, because $125\mid1000$. In either case $125\nmid N$.
In Lean the proof splits into the cases $L=1$, $L=2$, $L\ge3$ and the five values of $x$; in each of the $15$ cases \decl{simp} unfolds \eqref{eq:rec}, leaving a linear equation for $N$, and \decl{omega}, a decision procedure for linear integer arithmetic with divisibility by numerals, refutes $125\mid N$.
\end{proof}

For $q\ge0$ let
\[
  S_q=G(5,5q),
\]
which is \texttt{S q} in Lean. Its digit list is $\gamma(5,5q)=(5,7,9,1,3,5,7,9,1,3,\dots)$, so $S_q$ is the block $31975$ written $q$ times ($S_0=0$, $S_1=31975$, $S_2=3197531975$).

\begin{lemma}[\decl{S_succ}, \decl{dvd_S_25}, \decl{S_geom}]\label{lem:S}
For every $q\ge0$:
\begin{enumerate}
\item $S_{q+1}=31975+10^5S_q$;
\item $25\mid S_q$;
\item $99999\,S_q+31975=31975\cdot100000^q$.
\end{enumerate}
Moreover $S_q$ is admissible for $q\ge1$.
\end{lemma}

\begin{proof}
(1) Since $(3+2)\bmod10=5$, applying \eqref{eq:rec} five times gives $\gamma(5,5q+5)=(5,7,9,1,3,\gamma(5,5q))$ and $S_{q+1}=5+70+900+1000+30000+10^5S_q=31975+10^5S_q$.
(2) By induction on $q$: $S_0=0$, and $31975=25\cdot1279$, so $25\mid S_q$ implies $25\mid S_{q+1}$ by (1).
(3) By induction on $q$: for $q=0$ both sides equal $31975$. If (3) holds for $q$, then by (1)
\[
  99999\,S_{q+1}+31975=10^5\bigl(31975+99999\,S_q\bigr)=10^5\cdot31975\cdot100000^q=31975\cdot100000^{q+1}.
\]
Finally, $S_q=G(5,5q)$ with $5q\ge1$ is admissible by Lemma~\ref{lem:gen}.
\end{proof}

\begin{lemma}[\decl{exists_dvd_S}]\label{lem:euler}
Let $m'\ge1$ be coprime to $10$ and $q=\varphi(99999\,m')$, where $\varphi$ is Euler's totient function. Then $q\ge1$ and $25m'\mid S_q$.
\end{lemma}

\begin{proof}
Put $M=99999\,m'$. Since $99999=3^2\cdot41\cdot271$ and $m'$ are coprime to $10$, so is $M$; also $M\ge1$, so $q=\varphi(M)\ge1$. By Euler's theorem (\decl{Nat.ModEq.pow_totient}), $10^{q}\equiv1\pmod M$, hence $100000^q=(10^q)^5\equiv1\pmod M$. By Lemma~\ref{lem:S}(3),
\[
  99999\,S_q+31975=31975\cdot100000^q\equiv31975\pmod M,
\]
so $M\mid 99999\,S_q$, that is, $99999\,m'\mid99999\,S_q$, and therefore $m'\mid S_q$. Since $\gcd(25,m')=1$ and $25\mid S_q$ (Lemma~\ref{lem:S}(2)), we get $25m'\mid S_q$ (\decl{Nat.Coprime.mul_dvd_of_dvd_of_dvd}).
\end{proof}

In Lean, ``coprime to $10$'' is expressed by the two hypotheses \texttt{Nat.Coprime m' 2} and \texttt{Nat.Coprime m' 5}.

\begin{lemma}[\decl{exists_cofactor}]\label{lem:cof}
Let $m$ be odd with $125\nmid m$. Then there is an integer $m'$ coprime to $10$ with $m\mid 25m'$.
\end{lemma}

\begin{proof}
If $25\mid m$, write $m=25k$; then $k$ is odd, and $5\nmid k$ since otherwise $125\mid m$; take $m'=k$. If $5\mid m$ but $25\nmid m$, write $m=5k$; then $k$ is odd and $5\nmid k$, and $m\mid 25k$; take $m'=k$. Otherwise $m$ is coprime to $10$ and $m\mid 25m$; take $m'=m$.
\end{proof}

\begin{proof}[Proof of Theorem~\ref{thm:main}]
Let $n\ge0$ and $m=2n+1$. By Lemma~\ref{lem:zero} we must show that no admissible number is divisible by $m$ if and only if $125\mid m$.
If $125\mid m$ and $N$ is admissible with $m\mid N$, then $125\mid N$, contradicting Lemma~\ref{lem:125}.
If $125\nmid m$, then, as $m$ is odd, Lemma~\ref{lem:cof} gives $m'$ coprime to $10$ with $m\mid 25m'$, and Lemma~\ref{lem:euler} gives $q\ge1$ with $25m'\mid S_q$. Hence $m\mid S_q$, and $S_q$ is admissible by Lemma~\ref{lem:S}.
\end{proof}

\begin{remark}[A062886; not formalized]\label{rem:886}
The same argument proves Johnston's observation on A062886 \cite{A062886} for all $n$. There the digits are odd and each is two \emph{more} (mod $10$) than the previous one; as in Lemma~\ref{lem:gen}, the digit lists are $\bigl((x-2i)\bmod10\bigr)_{0\le i<L}$ with $x$ an odd digit and $L\ge1$. Such a number is either a digit or congruent modulo $100$ to its last two digits $13$, $35$, $57$, $79$ or $91$, so it is not divisible by $25$. Conversely, if $m=2n+1$ and $25\nmid m$, then $m\mid 5m'$ with $m'$ coprime to $10$, and the number $T_q$ formed by $q$ copies of the block $79135$ has the digit property, is divisible by $5$, and satisfies $99999\,T_q=79135\,(10^{5q}-1)$; as in Lemma~\ref{lem:euler}, $q=\varphi(99999\,m')$ gives $5m'\mid T_q$. Hence the zeros of A062886 are exactly the $n$ with $25\mid 2n+1$.
This remark is not part of the Lean development. We checked it, and Theorem~\ref{thm:main}, numerically for all $n<5000$ by a search over pairs (residue modulo $2n+1$, last digit).
\end{remark}

\section{Machine verification}\label{sec:lean}

\subsection*{Files and declarations}
The proof is in the file \path{lean/Math0/Frontier/OEIS/A062887.lean} of the math0 repository, \url{https://github.com/Foundation-0/math0}. It imports the frozen statement file \path{lean/Math0/Statements/OEIS/A062887.lean} (Section~\ref{sec:statement}) and the audit commands in \path{lean/Math0/Audit.lean}. The main declaration is \decl{Math0.Frontier.A062887.result}, of type \decl{Math0.Statements.A062887.statement}; the auxiliary declarations of Lemmas~\ref{lem:zero}--\ref{lem:cof} are private. The main theorem and the end of the file read, verbatim (docstring and blank lines omitted):
\begin{lstlisting}
theorem result : statement := by
  intro n
  rw [a_eq_zero_iff]
  constructor
  · intro hne
    by_contra h125
    obtain ⟨m', h2, h5, hdvd⟩ := exists_cofactor (m := 2 * n + 1) (by omega) h125
    obtain ⟨q, hq, hS⟩ := exists_dvd_S h2 h5
    exact hne ⟨S q, isDesc2_ofDigits_gen 5 (5 * q) (by norm_num) (by decide) (by omega),
      dvd_trans hdvd hS⟩
  · rintro h125 ⟨N, hN, hdvd⟩
    exact not_dvd_125 hN (dvd_trans h125 hdvd)
end Math0.Frontier.A062887
#math0_settles Math0.Frontier.A062887.result Math0.Statements.A062887.statement
\end{lstlisting}
The proof file uses no \texttt{native\_decide}, \texttt{sorry} or \texttt{set\_option}; its only uses of \texttt{decide} check that $5$ is odd and that the odd numbers below $10$ are $1,3,5,7,9$. The data check \decl{a_data} in the statement file uses \texttt{decide +kernel}, which evaluates in the kernel and adds no axioms.

\subsection*{Audit}
The command \texttt{\#math0\_settles}, defined in \path{lean/Math0/Audit.lean} for statements that math0 formalizes itself, fails unless (1) \decl{statement} is a \texttt{Prop}-valued definition without universe parameters; (2) \decl{statement} and \decl{result} use no axioms besides \decl{propext}, \decl{Classical.choice} and \decl{Quot.sound} (this rejects \decl{sorryAx}, \decl{Lean.ofReduceBool} from \texttt{native\_decide}, and custom axioms); and (3) the type of \decl{result} is syntactically the statement or its negation, up to metadata annotations and without unfolding. Here it reports
\begin{lstlisting}[basicstyle=\ttfamily\tiny]
math0 audit: 'Math0.Statements.A062887.statement' depends on axioms [Classical.choice, Quot.sound, propext]
math0 audit: 'Math0.Frontier.A062887.result' depends on axioms [Classical.choice, Quot.sound, propext]
math0 statement check: 'Math0.Frontier.A062887.result' proves 'Math0.Statements.A062887.statement'
\end{lstlisting}

\subsection*{Toolchain and reproduction}
We use Lean~4 \cite{Lean4} \texttt{v4.33.1} and Mathlib \cite{Mathlib} \texttt{v4.33.1}, pinned in \path{lean/}; the repository also pins formal-conjectures \cite{FC} at commit \texttt{c77fdd4}, which this result does not use.
To reproduce, run \path{scripts/verify.sh} from the repository root: it fetches the Mathlib build cache, builds all math0 results with \texttt{lake build} (re-running every audit command), checks that the pinned \texttt{formal\_conjectures} and \texttt{mathlib} checkouts are unmodified, and reports the number of passed statement checks. Each of the two files can also be checked with \texttt{lake env lean} in the directory \path{lean}, in well under a minute on our machine.

{\small\raggedright
\begin{thebibliography}{99}

\bibitem{FC}
M.~Firsching, P.~Lezeau, S.~Mercuri, M.~Z. Horv\'ath, Y.~Dillies, C.~S\"onne, E.~Wieser, F.~Zhang, T.~Hubert, B.~Ag\"uera y Arcas, and P.~Kohli, \emph{Formal Conjectures: An Open and Evolving Benchmark for Verified Discovery in Mathematics}, arXiv:2605.13171 [cs.AI], 2026. Repository: \url{https://github.com/google-deepmind/formal-conjectures}.

\bibitem{APN}
Google DeepMind, \emph{AlphaProof Nexus results}, \url{https://github.com/google-deepmind/alphaproof-nexus-results}.

\bibitem{Mathlib}
The mathlib Community, \emph{The Lean mathematical library}, in: Proceedings of the 9th ACM SIGPLAN International Conference on Certified Programs and Proofs (CPP 2020), ACM, 2020, 367--381.

\bibitem{Lean4}
L.~de Moura and S.~Ullrich, \emph{The Lean 4 theorem prover and programming language}, in: Automated Deduction -- CADE 28, Lecture Notes in Comput. Sci. \textbf{12699}, Springer, 2021, 625--635.

\bibitem{A062885}
The OEIS Foundation Inc., \url{https://oeis.org/A062885}, entry by A.~Murthy (2001); argument for $n=25$ by J.~W. Layman (2001); analysis file by D.~Reble.

\bibitem{A062886}
The OEIS Foundation Inc., \url{https://oeis.org/A062886}, entry by A.~Murthy (2001); comment and terms by N.~Johnston (May~19, 2011).

\bibitem{A062887}
The OEIS Foundation Inc., \url{https://oeis.org/A062887}, entry by A.~Murthy (June~28, 2001); more terms by S.~Kurz (2002); conjecture by S.~A. Irvine (April~14, 2023); revision \#12.

\end{thebibliography}
}

\end{document}
